% +--------------------------------------------------------------------+
% | Acknowledgements Page                                              |
% +--------------------------------------------------------------------+

\chapter*{Agradecimientos}

Este Trabajo de Fin de Grado no habría sido posible sin la ayuda y el apoyo de muchas personas
a las que queremos dedicar estas líneas.

En primer lugar, agradecemos a nuestro tutor el haber confiado en nuestra propuesta y haber
sabido darnos la libertad necesaria para explorar las decisiones de diseño que han ido
configurando este trabajo. Sus orientaciones en los momentos clave del proyecto, la paciencia
con las primeras versiones y el rigor con el que ha leído cada iteración de la memoria han
marcado el resultado final más de lo que cabe agradecer en un párrafo.

Agradecemos también al conjunto de profesores de la Facultad de Informática de la Universidad
Complutense de Madrid que, a lo largo de los cuatro años de la carrera, nos han transmitido
las bases sobre las que se sostiene este proyecto. La ingeniería del software, las bases de
datos, la teoría de la información y el análisis y diseño de interfaces aparecen aplicadas en
cada uno de los capítulos que siguen, fruto directo de su trabajo docente.

A nuestras familias les agradecemos el apoyo incondicional durante estos años, la paciencia en
los periodos de mayor carga y la confianza en que el camino emprendido valía la pena. Sin ese
soporte sostenido, llegar hasta este momento habría sido imposible.

A los compañeros y amigos que nos han acompañado en estos años de carrera, y muy especialmente
a quienes han sufrido nuestras explicaciones improvisadas sobre filtros facetados y navegación
hipermedia con genuino interés, gracias por hacer del recorrido algo más ligero.

Por último, queremos agradecer a las cinco personas que participaron desinteresadamente en el
estudio de usabilidad descrito en el capítulo de Pruebas y Validación. Sus observaciones,
preguntas y críticas dieron al trabajo una segunda vida y son el origen directo de varias de
las mejoras documentadas en esta memoria.
